\documentclass[../../main.tex]{subfiles}

\begin{document}

\chapter{Introduction}

The theoretical description of strongly correlated electron systems remains one of the most challenging and fascinating problems in condensed matter physics. Unlike weakly interacting systems, where electrons behave as nearly independent quasiparticles, strong Coulomb interactions lead to the emergence of collective phenomena such as the Mott metal–insulator transition \cite{Imada1998}, non-Fermi-liquid behavior, unconventional magnetism \cite{Liu2025} and high-temperature superconductivity \cite{Keimer2015}. Understanding these effects requires theoretical frameworks that go beyond single-particle approximations and simple mean-field descriptions.
\newline

Among modern approaches, dynamical mean-field theory (DMFT) \cite{Georges1996} has established itself as a cornerstone for modeling local electronic correlations. By mapping a lattice problem onto a quantum impurity model \cite{Georges1992}, DMFT provides non-perturbative access to single-particle observables, including spectral functions and quasiparticle renormalizations. However, DMFT inherently neglects non-local correlations, which are crucial for describing phenomena dominated by collective fluctuations, examples being pseudogap physics, spin-density waves, charge ordering and unconventional superconductivity.
\newline

To overcome this limitation, diagrammatic extensions of DMFT have been developed, among which the dynamical vertex approximation ($\dga$) \cite{Toschi2007, Rohringer2018} plays a central role. $\dga$ systematically incorporates non-local correlations by taking the local two-particle irreducible vertex from DMFT as a building block and constructing ladder or parquet diagrammatic expansions to compute momentum- and frequency-resolved quantities. In this way, $\dga$ grants access to susceptibilities, self-energies and pairing instabilities beyond the local approximation, enabling the study of competing phases on equal footing \cite{Rohringer2016}.
\newline

While $\dga$ has been successfully applied to single-orbital systems \cite{Worm2024, DiCataldo2024}, many strongly correlated materials exhibit a pronounced multi-orbital character. This is particularly relevant for realistic compounds such as the recently discovered (infinite-layer) nickelates \cite{Christiansson2023, Li2019}, where orbital selectivity, Hund’s coupling and intertwined instabilities strongly influence low-energy physics and superconducting pairing mechanisms \cite{DeMedici2011}. Accurately describing these systems requires a multi-orbital generalization of $\dga$, which significantly increases computational and algorithmic complexity due to the higher dimensionality of the involved vertex functions.
\newline

This thesis presents the development of a self-consistent multi-orbital ladder-$\dga$ framework applied to the multi-orbital Hubbard model. The implementation focuses on numerical efficiency, scalability and flexibility to enable calculations at very low temperatures and for realistic multi-orbital setups.
\newline

This work is structured as follows: \chref{ch:foundations_and_models} introduces the theoretical background of condensed matter models and the necessary foundation for describing correlated electron systems. We mainly focus on the multi-orbital Hubbard Hamiltonian and its various simplifications, including the well-known Kanamori parametrization. We then develop the theoretical tools required for realistic $\dga$ calculations, including one- and two-particle Green’s functions, susceptibilities, the Bethe-Salpeter equation and the Schwinger-Dyson equation in \chref{ch:quantum_many_body_framework}. We also cover vertex asymptotics and explain their role in reducing numerical complexity while simultaneously significantly enhancing the capabilities of low-temperature calculations. Finally, the ladder $\dga$ formalism and its connection to superconducting pairing instabilities via the Eliashberg equation are derived. Afterwards, in \chref{ch:the_code}, we describe the design and implementation of the computational framework which was developed in the course of this thesis. We explain the parallelization strategy using MPI, detail the data structures and APIs and discuss optimization techniques such as symmetry exploitation and vertex asymptotics. This chapter also presents benchmark results, scaling analyses and a validation of the implementation against known reference cases. Lastly, we summarize the main contributions and outline possible future extensions in \chref{ch:conclusions_and_outlook}.

\end{document}